\section{Riformulazione: entropia stabilizer e norme di Gowers}\label{sec:sre}

Per $\psi\in\C^{2^n}$ sia $f_\psi:\Fb_2^n\to\C$, $f_\psi(x)=2^{n/2}\braket{x}{\psi}$, così che $\E_x|f_\psi(x)|^2=1$. Indichiamo con $\mathcal C$ la coniugazione complessa. La norma di Gowers $U^3$ è
\[
\|f\|_{U^3}^8=\E_{x,h_1,h_2,h_3\in\Fb_2^n}\ \prod_{\omega\in\{0,1\}^3}\mathcal C^{|\omega|}f\bigl(x+\omega_1h_1+\omega_2h_2+\omega_3h_3\bigr).
\]

\begin{proposition}[SRE $=$ Gowers $U^3$]\label{prop:gowers}
Per ogni stato normalizzato $\psi$,
\[
2^{-M_2(\psi)}=\|f_\psi\|_{U^3}^8,\qquad\text{cioè}\qquad M_2(\psi)=-8\log_2\|f_\psi\|_{U^3}.
\]
\end{proposition}

\noindent La quantità $\|f_\psi\|_{U^3}$ è la ``norma di Gowers di uno stato quantistico'' di \cite{AD24}. Lì il legame con il quarto momento dello spettro di Pauli è implicito; qui lo rendiamo esplicito.

\begin{proof}
A meno di una fase, $\bra\psi X^aZ^b\ket\psi=\hat g_a(b)$, dove $g_a(x)=\overline{\psi(x+a)}\,\psi(x)$ e $\hat g(b)=\sum_xg(x)(-1)^{b\cdot x}$. Per ogni $g$ vale l'identità
\[
\sum_b|\hat g(b)|^4=2^n\sum_{x,h_1,h_2}g(x)\,\overline{g(x+h_1)}\,\overline{g(x+h_2)}\,g(x+h_1+h_2).
\]
Inoltre $g_a=2^{-n}\Delta_af$, con $f=f_\psi$ e $\Delta_af(x)=f(x)\overline{f(x+a)}$. Sostituendo,
\[
2^{-n}\sum_{a,b}|\hat g_a(b)|^4=\E_a\|\Delta_af\|_{U^2}^4=\|f\|_{U^3}^8 .
\]
Il primo membro è $2^{-n}\sum_P\bra\psi P\ket\psi^4=2^{-M_2(\psi)}$.
\end{proof}

\begin{proposition}\label{prop:M2F}
Per ogni stato normalizzato $\psi$,
\[
\tfrac14M_2(\psi)\ \le\ \log_2\frac1{F(\psi)}\ \le\ \log_2\xi(\psi).
\]
\end{proposition}

\begin{proof}
Sia $s$ uno stato stabilizer che realizza $F$, e sia $S$ il suo gruppo stabilizzatore con segni, formato da $2^n$ operatori di Pauli con $g\ket s=\ket s$. Allora $\ket s\bra s=2^{-n}\sum_{g\in S}g$, e quindi $F=2^{-n}\sum_{g\in S}\bra\psi g\ket\psi$. Poiché $x\mapsto x^4$ è convessa, la disuguaglianza di Jensen dà
\[
2^{-n}\sum_{P}\bra\psi P\ket\psi^4\ \ge\ 2^{-n}\sum_{g\in S}\bra\psi g\ket\psi^4\ \ge\ \Bigl(2^{-n}\sum_{g\in S}\bra\psi g\ket\psi\Bigr)^4=F^4.
\]
La seconda disuguaglianza dell'enunciato è \ref{R2}.
\end{proof}

\noindent Vale anche una relazione inversa polinomiale. Il teorema inverso per la norma di Gowers $U^3$ degli stati quantistici \cite{AD24}, che passa per la congettura di Marton (PFR) su $\Fb_2^n$ \cite{GGMT23}, dà $F(\psi)\ge\poly\bigl(2^{-M_2(\psi)}\bigr)$. Quindi $\log_2(1/F)$ e $M_2$ coincidono a meno di costanti moltiplicative e additive. Le costanti in \cite{AD24} sono grandi; qui conta solo che la relazione sia polinomiale.

\subsection{Tre congetture equivalenti a meno di costanti}

\begin{conjecture}\label{conj:main}
Esistono costanti $C,D$ tali che, per ogni $n$ e ogni stato puro $\psi$ su $n$ qubit:
\begin{enumerate}[(A),leftmargin=2em]
\item\label{cA} (SRE--rango) $M_2(\psi)\le C\log_2\chi(\psi)+D$;
\item\label{cB} (fidelity--rango) $\log_2(1/F(\psi))\le C\log_2\chi(\psi)+D$;
\item\label{cC} (extent--rango, domanda aperta in \cite{SL26}) $\log_2\xi(\psi)\le C\log_2\chi(\psi)+D$.
\end{enumerate}
\end{conjecture}

\noindent Per la \cref{prop:M2F} si ha \ref{cC}$\Rightarrow$\ref{cB}$\Rightarrow$\ref{cA}, e per il teorema inverso di \cite{AD24} vale anche \ref{cA}$\Rightarrow$\ref{cB}, con costanti diverse.

\begin{theorem}\label{thm:conj-exp}
Ciascuna delle tre congetture implica $\chi(\phi^{\otimes n})=2^{\Omega(n)}$ per ogni $\phi\notin\Stab_1$. Per $\Hst$, la \ref{cA} dà $\chi(\Hst^{\otimes n})\ge2^{(n\log_2(4/3)-D)/C}$ e la \ref{cB} dà $\chi(\Hst^{\otimes n})\ge2^{(0.2284\,n-D)/C}$.
\end{theorem}

\begin{proof}
$M_2$ è additiva e strettamente positiva su $\phi\notin\Stab_1$. Fidelity ed extent sono moltiplicativi sui prodotti di stati a un qubit \cite{BBCCGH19}. Si usano poi i valori \eqref{eq:values}.
\end{proof}

\paragraph{Cosa si sa.} Il risultato di Sinha--Lovitz, insieme alla \cref{prop:M2F}, dà $M_2(\psi)\le O\bigl(\sqrt{k\log k}\bigr)$ per stati di rango $k$. Kalra--Sinha danno $\log_2(1/F)\le O(k\log k)$ \cite{KS25}. Tutto il gap fra il lower bound quadratico e quello esponenziale sta nella differenza fra $\sqrt k$ e $\log k$.

\paragraph{Vincoli sulle costanti.}
\begin{itemize}[leftmargin=1.5em]
\item Per somme di $k$ stati stabilizer quasi ortogonali con pesi uguali, i contributi diagonali danno $2^{-M_2}\approx\sum_j|c_j|^8=k^{-3}$, cioè $M_2\approx3\log_2k$. Numericamente, a $n=12$, si trova $M_2=2.97,\ 4.74,\ 5.93$ per $k=2,3,4$ (\cref{app:num}). Quindi serve $C\ge3$ in \ref{cA}, e $C\ge1$ in \ref{cB}, perché in quel caso $F\approx1/k$.
\item \emph{Forma forte.} Si potrebbe congetturare $F(\psi)\,\chi(\psi)\ge1$, cioè \ref{cB} con $C=1$ e $D=0$. Darebbe $\chi(\Hst^{\otimes n})\ge(4-2\sqrt2)^n\approx1.1716^n$, che coincide con lo scaling del rango approssimato. I dati sono vicini al limite:
\begin{itemize}
\item per $\Hst$: $F^m\chi=1.71,\ 1.46,\ 2.12$ per $m=1,2,4$;
\item per lo stato faccia, con $F=\tfrac12(1+1/\sqrt3)\approx0.789$ e i valori di $\chi$ di \cite{LR26}: $F^m\chi=1.58,\ 1.24,\ 1.47,\ 1.16$ per $m=1,\dots,4$;
\item per lo stato $W_n$: $F(W_n)=9/(4n)$, mentre $\chi(W_n)\le\lceil n/2\rceil$ (somme di coppie), quindi $F\chi\le9/8+O(1/n)$.
\end{itemize}
Una decomposizione delle potenze dello stato faccia con esponente minore di $\log_2(1/F)\approx0.342$ per copia confuterebbe la forma forte, ma non \ref{cB}.
\end{itemize}

\subsection{Formulazione additivo-combinatoria}

\begin{conjecture}[forma Gowers]\label{conj:gowers}
Sia $f=\sum_{j=1}^kc_jg_j$, dove ogni $g_j$ è una funzione stabilizer normalizzata in $L^2$: una fase quadratica, eventualmente non classica a valori in $\Z_4$, ristretta a un sottospazio affine di $\Fb_2^n$. Se $\E|f|^2=1$, allora $\|f\|_{U^3}\ge k^{-O(1)}$.
\end{conjecture}

\noindent Per la \cref{prop:gowers} questa è esattamente la \ref{cA}. Alcune osservazioni.
\begin{itemize}[leftmargin=1.5em]
\item Il bound noto è $\|f\|_{U^3}^8\ge F^4\ge\xi^{-4}\ge2^{-O(\sqrt{k\log k})}$ \cite{SL26}.
\item L'analogo per $U^2$ con i caratteri è banale: se $f=\sum_jc_j\chi_j$ con caratteri ortonormali, allora $\|f\|_{U^2}^4=\sum_j|c_j|^4\ge1/k$.
\item La difficoltà sta in due punti. Primo, $f$ non è limitata ($|f|$ arriva a $2^{n/2}$), mentre i teoremi inversi standard sono per funzioni limitate. Secondo, le cancellazioni fra i $g_j$ devono essere esponenziali (\cref{prop:cond-number}).
\item Una risposta positiva dà $\chi(\Tst^{\otimes n})=2^{\Omega(n)}$. Un controesempio, cioè somme di $k$ funzioni stabilizer con norma $U^3$ pari a $k^{-\omega(1)}$, mostrerebbe che nessuna disuguaglianza generale fra rango e $M_2$ (o fidelity, o extent) può dare un bound esponenziale.
\end{itemize}
